\chapter{Complex abelian varieties without coordinates}\label{app:abelian}

The purpose of this appendix is to explain in which space every object of
\cref{sec:conventions} naturally lives. We start from an abstract complex
abelian variety, with no basis and no presentation as a quotient
$\CC^{g}/\Lambda$, and introduce coordinates only by choosing a
\emph{frame} of the first homology group. The main point is that the period
map $\Homology{A}{\ZZ} \to \LieAlg{A}$ is canonical, whereas a formula like
$v \mapsto \int_v dz$ already involves the choice of a differential
$dz$, i.e.\ of a coordinate on $\LieAlg{A}$. References for the standard
facts are \cite[\S\S 1--4]{Mumford1974} and \cite[Ch.~1, 2, 4]{BirkenhakeLange2004}.

\section{The canonical objects}

\begin{setup}[Complex abelian variety]\label{setup:abelian}
Let $A$ be a complex abelian variety of dimension $g \geq 1$. We use:
\begin{enumerate}
\item $\LieAlg{A}$, the tangent space of $A$ at the origin, a complex vector
  space of dimension $g$;
\item $\Differentials{A}$, the space of holomorphic $1$-forms on $A$;
\item $\Homology{A}{\ZZ} := H_1(A(\CC), \ZZ)$, a free $\ZZ$-module of rank
  $2g$, and $\Homology{A}{R} := \Homology{A}{\ZZ}\otimes_{\ZZ} R$ for
  $R \in \{\QQ, \RR, \CC, \ProfiniteIntegers, \FiniteAdeles{\QQ}\}$;
\item $\Cohomology{A}{R} := \Hom(\Homology{A}{\ZZ}, R)$;
\item $\ExpMap{A} \colon \LieAlg{A} \to A(\CC)$, the exponential map of the
  connected commutative complex Lie group $A(\CC)$.
\end{enumerate}
\end{setup}

\begin{lemma}[Differentials are the cotangent space]\label{lem:differentials-cotangent}
\slogan{A holomorphic $1$-form on an abelian variety is determined by its
value at the origin.}
In \cref{setup:abelian}, every $\omega \in \Differentials{A}$ is invariant
under translations, and $\omega \mapsto \omega(0)$ is an isomorphism
$\Differentials{A} \to \Hom_{\CC}(\LieAlg{A}, \CC)$.
\end{lemma}

\begin{proof}
\cite[\S 1]{Mumford1974}, \cite[\S 1.4]{BirkenhakeLange2004}.
\end{proof}

\begin{lemma}[Uniformisation]\label{lem:uniformisation}
\slogan{$A(\CC)$ is its Lie algebra modulo its first homology.}
In \cref{setup:abelian}, $\ExpMap{A}$ is a surjective holomorphic
homomorphism and a universal covering. For a loop $\gamma$ in $A(\CC)$ based
at $0$ let $\PeriodInclusion{A}([\gamma]) \in \LieAlg{A}$ be the endpoint of
the lift of $\gamma$ through $\ExpMap{A}$ starting at $0$. Then
$\PeriodInclusion{A}$ is a well defined injective homomorphism, the sequence
\[
  0 \to \Homology{A}{\ZZ} \xrightarrow{\ \PeriodInclusion{A}\ } \LieAlg{A}
  \xrightarrow{\ \ExpMap{A}\ } A(\CC) \to 0
\]
is exact, and the $\RR$-linear extension
$\Homology{A}{\RR} \to \LieAlg{A}$ of $\PeriodInclusion{A}$ is an isomorphism
of real vector spaces.
\end{lemma}

\begin{proof}
\cite[\S 1]{Mumford1974}: $\ExpMap{A}$ is the universal covering, its kernel
is the group of deck transformations, which is
$\pi_1(A(\CC), 0) = \Homology{A}{\ZZ}$ since $\pi_1$ is abelian, and the
kernel is a lattice because $A(\CC)$ is compact. Compare
\cite[IV.5.1]{DeligneRapoport1973}.
\end{proof}

\begin{proposition}[The period map is canonical]\label{prop:canonical-period-map}
\slogan{Integrating differentials over cycles is the inclusion of the
homology into the Lie algebra; no choice of differential is needed.}
In \cref{setup:abelian}, for $v \in \Homology{A}{\ZZ}$ and
$\omega \in \Differentials{A}$ one has
$\Period{v}{\omega} = \omega(0)\bigl(\PeriodInclusion{A}(v)\bigr)$. Hence,
under the identification $\LieAlg{A} = \Hom_{\CC}(\Differentials{A}, \CC)$
given by the cotangent space description of differentials
(\cref{lem:differentials-cotangent}),
\[
  \PeriodInclusion{A}(v) = \bigl(\omega \mapsto \Period{v}{\omega}\bigr).
\]
If $g = 1$ and $\omega \in \Differentials{A}$ is nonzero, then $\omega(0)$ is
a linear coordinate $z$ on $\LieAlg{A}$ with $\ExpMap{A}^{*}\omega = dz$,
and $v \mapsto \Period{v}{\omega}$ identifies $\Homology{A}{\ZZ}$ with a
lattice $\Lambda_{\omega} \subseteq \CC$ such that $A(\CC) \cong \CC/\Lambda_{\omega}$.
\end{proposition}

\begin{proof}
Since $\ExpMap{A}$ is a homomorphism and $\omega$ is translation invariant,
$\ExpMap{A}^{*}\omega$ is the translation invariant $1$-form on the vector
space $\LieAlg{A}$ with value $\omega(0)$ at $0$, i.e.\ the constant form
$\omega(0)$. If $\tilde\gamma$ is the lift of a loop $\gamma$ starting at
$0$, then
$\Period{\gamma}{\omega} = \Period{\tilde\gamma}{\ExpMap{A}^{*}\omega} =
\omega(0)(\PeriodInclusion{A}([\gamma]) - 0)$. The last statement follows
from uniformisation (\cref{lem:uniformisation}).
\end{proof}

\begin{remark}
So writing an elliptic curve as $\CC/\Lambda$ and integrating $dz$ amounts
to choosing a nonzero element of the one-dimensional space
$\Differentials{A}$; the canonical object is the map
$\PeriodInclusion{A} \colon \Homology{A}{\ZZ} \to \LieAlg{A}$.
\end{remark}

\section{The Hodge filtration}

\begin{definition}[Complex structure and Hodge decomposition]\label{def:hodge}
In \cref{setup:abelian}, let $\ComplexStructure{A}$ be the $\RR$-linear
automorphism of $\Homology{A}{\RR}$ corresponding to multiplication by $i$
on $\LieAlg{A}$ under the isomorphism of \cref{lem:uniformisation}. On
$\Homology{A}{\CC} = \Homology{A}{\RR}\otimes_{\RR}\CC$ let
$\HodgePiece{A}{-1}{0}$ and $\HodgePiece{A}{0}{-1}$ be the eigenspaces of
$\ComplexStructure{A}\otimes\CC$ for the eigenvalues $i$ and $-i$. The
\emph{Hodge filtration} is $\HodgeFiltration{A} := \HodgePiece{A}{0}{-1}$.
Let $\Complexification{\PeriodInclusion{A}} \colon \Homology{A}{\CC} \to
\LieAlg{A}$ be the $\CC$-linear extension of $\PeriodInclusion{A}$, and
$\HodgeOneZero{A} \subseteq \Cohomology{A}{\CC}$ the set of functionals
$v \mapsto \Period{v}{\omega}$ with $\omega \in \Differentials{A}$.
\end{definition}

\begin{lemma}[The Hodge filtration is the kernel of the period map]
\label{lem:hodge-kernel}
\slogan{$\LieAlg{A} = \Homology{A}{\CC}/F^{0}$ and the holomorphic
differentials are the functionals vanishing on $F^{0}$.}
In \cref{setup:abelian},
$\Homology{A}{\CC} = \HodgePiece{A}{-1}{0}\oplus\HodgePiece{A}{0}{-1}$, both
summands have dimension $g$ and are interchanged by complex conjugation,
$\HodgeFiltration{A} = \{x + i\ComplexStructure{A}x : x \in \Homology{A}{\RR}\}$,
the map $\Complexification{\PeriodInclusion{A}}$ is surjective with kernel
$\HodgeFiltration{A}$, and $\HodgeOneZero{A}$ is the annihilator of
$\HodgeFiltration{A}$ in $\Cohomology{A}{\CC}$. In particular
$\HodgeFiltration{A} \cap \Homology{A}{\RR} = 0$.
\end{lemma}

\begin{proof}
$\ComplexStructure{A}^{2} = -1$, so the eigenspace decomposition exists,
and $\ComplexStructure{A}(x + i\ComplexStructure{A}x) =
-i(x + i\ComplexStructure{A}x)$ shows the description of
$\HodgeFiltration{A}$ and the dimension count. By definition of
$\ComplexStructure{A}$ one has $\Complexification{\PeriodInclusion{A}}\circ
(\ComplexStructure{A}\otimes\CC) = i\,\Complexification{\PeriodInclusion{A}}$;
on the $(-i)$-eigenspace this forces
$\Complexification{\PeriodInclusion{A}} = 0$. The map is surjective since
its restriction to $\Homology{A}{\RR}$ is, so its kernel has dimension $g$
and equals $\HodgeFiltration{A}$. By the canonical period map
(\cref{prop:canonical-period-map}) $\HodgeOneZero{A}$ consists of the
functionals $\lambda\circ\Complexification{\PeriodInclusion{A}}$ with
$\lambda \in \Hom_{\CC}(\LieAlg{A}, \CC)$, i.e.\ of the functionals
vanishing on the kernel. A real vector $x$ with
$\ComplexStructure{A}x = -ix$ is zero.
\end{proof}

\begin{remark}
\todo[inline]{Check that the labelling of $\HodgePiece{A}{-1}{0}$ and
$\HodgePiece{A}{0}{-1}$ agrees with Deligne's conventions for the Hodge
structure on $H_1$ in \cite{DeligneVarietes1979}.}
\end{remark}

\section{Riemann forms}

\begin{definition}[Riemann form]\label{def:riemann-form}
In \cref{setup:abelian}, a \emph{Riemann form} on $A$ is an alternating
bilinear form
$\RiemannForm{} \colon \Homology{A}{\ZZ}\times\Homology{A}{\ZZ} \to \ZZ$ such
that, for its $\RR$-bilinear extension,
$\RiemannForm{}(\ComplexStructure{A}x, \ComplexStructure{A}y) =
\RiemannForm{}(x,y)$ for all $x, y$ and
$\RiemannForm{}(\ComplexStructure{A}x, x) > 0$ for all $x \neq 0$. It is
\emph{principal} if it is unimodular. We write $\RiemannForm{}$ also for
its $\CC$-bilinear extension to $\Homology{A}{\CC}$.
\end{definition}

\begin{remark}[Sign convention]\label{rem:riemann-sign}
This is the sign convention of \cite[\S 2]{Mumford1974} and
\cite[\S 2.1]{BirkenhakeLange2004}: the form
$H(x,y) := \RiemannForm{}(\ComplexStructure{A}x, y) + i\RiemannForm{}(x,y)$
on $\LieAlg{A}$ is hermitian (complex linear in $x$) and positive definite.
For an elliptic curve and a basis $a, b$ of $\Homology{A}{\ZZ}$ with
intersection number $\IntersectionPairing{a}{b} = 1$ for the complex
orientation, one finds $\RiemannForm{}(a,b) = -1$ for the principal Riemann
form of \cref{lem:elliptic-riemann-form}: the principal Riemann form is
\emph{minus} the intersection pairing.
\todo[inline]{Double check the sign relative to the intersection pairing
against a standard reference.}
\end{remark}

\begin{lemma}[Riemann relations]\label{lem:riemann-relations}
\slogan{The Hodge filtration is a positive Lagrangian.}
In \cref{setup:abelian}, let $\RiemannForm{}$ be a Riemann form on $A$.
Then $\HodgeFiltration{A}$ is a Lagrangian subspace of $\Homology{A}{\CC}$
for $\RiemannForm{}$, and $-i\,\RiemannForm{}(u, \overline{u}) > 0$ for every
nonzero $u \in \HodgeFiltration{A}$.
\end{lemma}

\begin{proof}
Write $u = x + i\ComplexStructure{A}x$ and $u' = y + i\ComplexStructure{A}y$
(\cref{lem:hodge-kernel}). Using the invariance of $\RiemannForm{}$ under
$\ComplexStructure{A}$ and $\ComplexStructure{A}^{2} = -1$ one gets
$\RiemannForm{}(u, u') = i(\RiemannForm{}(x, \ComplexStructure{A}y) +
\RiemannForm{}(\ComplexStructure{A}x, y)) = 0$ and
$\RiemannForm{}(u, \overline{u}) = 2i\,\RiemannForm{}(\ComplexStructure{A}x, x)$.
Isotropic of dimension $g$ in dimension $2g$ means Lagrangian.
\end{proof}

\begin{proposition}[Poincar\'e duality for the Hodge filtration]
\label{prop:poincare-hodge}
\slogan{Given a Riemann form, a holomorphic differential is the same thing
as a vector in the Hodge filtration of $H_1$.}
In \cref{setup:abelian}, let $\RiemannForm{}$ be a Riemann form on $A$. For
every $\omega \in \Differentials{A}$ there is a unique
$\PoincareDual{\RiemannForm{}}{\omega} \in \HodgeFiltration{A}$ with
\[
  \Period{v}{\omega} = \RiemannForm{}\bigl(\PoincareDual{\RiemannForm{}}{\omega}, v\bigr)
  \qquad\text{for all } v \in \Homology{A}{\ZZ},
\]
and $\omega \mapsto \PoincareDual{\RiemannForm{}}{\omega}$ is a $\CC$-linear
isomorphism $\Differentials{A} \to \HodgeFiltration{A}$.
\end{proposition}

\begin{proof}
By positivity $\RiemannForm{}$ is nondegenerate, so
$\PoincareMap{\RiemannForm{}} \colon \Homology{A}{\CC} \to \Cohomology{A}{\CC}$,
$u \mapsto \RiemannForm{}(u, \cdot)$, is an isomorphism. By the Riemann
relations (\cref{lem:riemann-relations}) it maps $\HodgeFiltration{A}$ into
the annihilator of $\HodgeFiltration{A}$, which has the same dimension $g$
and equals $\HodgeOneZero{A}$ by \cref{lem:hodge-kernel}. Finally
$\omega \mapsto (v \mapsto \Period{v}{\omega})$ is an isomorphism
$\Differentials{A} \to \HodgeOneZero{A}$ by the same lemma.
\end{proof}

\begin{lemma}[The Riemann form of an elliptic curve]\label{lem:elliptic-riemann-form}
\slogan{An elliptic curve has a canonical principal Riemann form.}
In \cref{setup:abelian} with $g = 1$ there is a unique principal Riemann
form $\RiemannForm{A}$. For every $\CC$-linear isomorphism
$z \colon \LieAlg{A} \to \CC$, with $\Lambda := z(\PeriodInclusion{A}(\Homology{A}{\ZZ}))$,
\[
  \RiemannForm{A}(x, y) = \frac{\mathrm{Im}\bigl(z(\PeriodInclusion{A}x)\,
  \overline{z(\PeriodInclusion{A}y)}\bigr)}{\Covolume{\Lambda}} .
\]
\end{lemma}

\begin{proof}
The alternating forms on a free module of rank $2$ are the multiples of a
generator of $\Hom(\textstyle\bigwedge^{2}\Homology{A}{\ZZ}, \ZZ) \cong \ZZ$, so
the unimodular ones are $\pm$ a generator. Invariance under
$\ComplexStructure{A}$ is automatic since $\det\ComplexStructure{A} = 1$, and
positivity fixes the sign. The displayed form is independent of $z$
(rescaling $z$ by $\lambda$ multiplies numerator and covolume by
$|\lambda|^{2}$), alternating, $\ComplexStructure{A}$-invariant, and
satisfies $\RiemannForm{A}(\ComplexStructure{A}x, x) =
|z(\PeriodInclusion{A}x)|^{2}/\Covolume{\Lambda} > 0$; for a basis
$\lambda_1, \lambda_2$ of $\Lambda$ one has
$|\mathrm{Im}(\lambda_1\overline{\lambda_2})| = \Covolume{\Lambda}$, so it is
unimodular.
\end{proof}

\section{Torsion points and Tate modules}

We first describe the construction abstractly.

\begin{setup}[Free module]\label{setup:free-module}
Let $\FreeModule$ be a free $\ZZ$-module of finite rank. Put
\[
  \DivisibleGroup{\FreeModule} := (\FreeModule\otimes\QQ)/\FreeModule,
  \qquad
  \DivTorsion{\FreeModule}{N} := \{x \in \DivisibleGroup{\FreeModule} : Nx = 0\}
  \quad (N \geq 1),
\]
and let $\AbstractTate{\FreeModule} := \varprojlim_N \DivTorsion{\FreeModule}{N}$
and $\pAdicTate{p}{\FreeModule} := \varprojlim_k \DivTorsion{\FreeModule}{p^{k}}$
($p$ prime) be the Tate modules, the transition maps
$\DivTorsion{\FreeModule}{NM} \to \DivTorsion{\FreeModule}{N}$ being
multiplication by $M$. For $N \geq 1$ let
$\DivisionMap{N} \colon \FreeModule/N\FreeModule \to
\tfrac{1}{N}\FreeModule/\FreeModule$ be division by $N$,
$v \bmod N\FreeModule \mapsto v/N \bmod \FreeModule$.
\end{setup}

\begin{lemma}[Torsion of $(\FreeModule\otimes\QQ)/\FreeModule$]\label{lem:divisible-torsion}
\slogan{The $N$-torsion of $(\mathcal{H}\otimes\QQ)/\mathcal{H}$ is
$\frac{1}{N}\mathcal{H}/\mathcal{H}$.}
In \cref{setup:free-module}, $\DivisibleGroup{\FreeModule}$ is a divisible
torsion group,
$\DivTorsion{\FreeModule}{N} = \tfrac{1}{N}\FreeModule/\FreeModule$, the
inclusion $\DivTorsion{\FreeModule}{N} \subseteq \DivTorsion{\FreeModule}{NM}$
is the inclusion of subgroups, multiplication by $M$ maps
$\DivTorsion{\FreeModule}{NM}$ onto $\DivTorsion{\FreeModule}{N}$, and
$\DivisibleGroup{\FreeModule} = \bigcup_N \DivTorsion{\FreeModule}{N}
\cong \FreeModule\otimes\QQ/\ZZ$.
\end{lemma}

\begin{proof}
$x \in \FreeModule\otimes\QQ$ satisfies $Nx \in \FreeModule$ if and only if
$x \in \tfrac{1}{N}\FreeModule$. Divisibility and the last isomorphism hold
because $\FreeModule \cong \ZZ^{r}$ and $\QQ/\ZZ$ is divisible.
\end{proof}

\begin{lemma}[Division by $N$]\label{lem:division-by-N}
\slogan{Division by powers of $p$ identifies the tower of reductions of
$\mathcal{H}$ with the tower of $p$-power torsion of
$(\mathcal{H}\otimes\QQ)/\mathcal{H}$ under multiplication by $p$.}
In \cref{setup:free-module}, $\DivisionMap{N}$ is an isomorphism for every
$N \geq 1$. For $N, M \geq 1$ it transforms reduction
$\FreeModule/NM\FreeModule \to \FreeModule/N\FreeModule$ into
multiplication by $M$ from $\DivTorsion{\FreeModule}{NM}$ to
$\DivTorsion{\FreeModule}{N}$, and multiplication by $M$ from
$\FreeModule/N\FreeModule$ to $\FreeModule/NM\FreeModule$ into the inclusion
$\DivTorsion{\FreeModule}{N} \subseteq \DivTorsion{\FreeModule}{NM}$. In
particular, for a prime $p$ the following diagram commutes and its
vertical maps are isomorphisms:
\[
\begin{tikzcd}[column sep=large]
  \cdots \arrow[r, "\cdot p"] &
  \tfrac{1}{p^{3}}\FreeModule/\FreeModule \arrow[r, "\cdot p"] &
  \tfrac{1}{p^{2}}\FreeModule/\FreeModule \arrow[r, "\cdot p"] &
  \tfrac{1}{p}\FreeModule/\FreeModule \arrow[r] & 0 \\
  \cdots \arrow[r, "\mathrm{red}"] &
  \FreeModule/p^{3}\FreeModule \arrow[r, "\mathrm{red}"]
    \arrow[u, "\DivisionMap{p^{3}}"', "\wr"] &
  \FreeModule/p^{2}\FreeModule \arrow[r, "\mathrm{red}"]
    \arrow[u, "\DivisionMap{p^{2}}"', "\wr"] &
  \FreeModule/p\FreeModule \arrow[r]
    \arrow[u, "\DivisionMap{p}"', "\wr"] & 0
\end{tikzcd}
\]
Here $\mathrm{red}$ denotes reduction and $\cdot p$ multiplication by $p$.
\end{lemma}

\begin{proof}
$\DivisionMap{N}$ is well defined and bijective since multiplication by
$1/N$ is a bijection $\FreeModule \to \tfrac{1}{N}\FreeModule$ mapping
$N\FreeModule$ onto $\FreeModule$. For $v \in \FreeModule$ one has
$M\cdot(v/NM) = v/N$, which is the first compatibility, and
$(Mv)/(NM) = v/N$, which is the second.
\end{proof}

\begin{corollary}[Tate modules of $(\FreeModule\otimes\QQ)/\FreeModule$]
\label{cor:abstract-tate}
\slogan{The Tate module of $(\mathcal{H}\otimes\QQ)/\mathcal{H}$ is
$\mathcal{H}\otimes\widehat{\ZZ}$.}
In \cref{setup:free-module}, the division maps (\cref{lem:division-by-N})
induce isomorphisms
$\FreeModule\otimes\ProfiniteIntegers = \varprojlim_N\FreeModule/N\FreeModule
\to \AbstractTate{\FreeModule}$ and
$\FreeModule\otimes\pAdicIntegers{p} \to \pAdicTate{p}{\FreeModule}$. They
extend to an isomorphism
$\FreeModule\otimes\FiniteAdeles{\QQ} \to \AbstractTate{\FreeModule}\otimes\QQ$,
and $(\FreeModule\otimes\FiniteAdeles{\QQ})/(\FreeModule\otimes\ProfiniteIntegers)
\cong \DivisibleGroup{\FreeModule}$.
\end{corollary}

\begin{proof}
The inverse limits agree by the compatibility of division with the
transition maps (\cref{lem:division-by-N}); $\varprojlim_N \FreeModule/N\FreeModule
= \FreeModule\otimes\ProfiniteIntegers$ because $\FreeModule$ is free of finite
rank. Tensor with $\QQ$, and use
$\FiniteAdeles{\QQ}/\ProfiniteIntegers \cong \QQ/\ZZ$ and the torsion of
$\DivisibleGroup{\FreeModule}$ (\cref{lem:divisible-torsion}).
\end{proof}

We now apply this to $\FreeModule = \Homology{A}{\ZZ}$.

\begin{lemma}[Torsion points are rational homology modulo integral homology]
\label{lem:torsion-points}
\slogan{$A(\CC)_{\mathrm{tors}} = H_1(A,\QQ)/H_1(A,\ZZ)$ and
$A[N] = \tfrac{1}{N}H_1(A,\ZZ)/H_1(A,\ZZ)$.}
In \cref{setup:abelian}, the $\QQ$-linear extension of $\PeriodInclusion{A}$
followed by $\ExpMap{A}$ induces an isomorphism
$\DivisibleGroup{\Homology{A}{\ZZ}} = \Homology{A}{\QQ}/\Homology{A}{\ZZ}
\to A(\CC)_{\mathrm{tors}}$ of groups, which restricts to isomorphisms
$\tfrac{1}{N}\Homology{A}{\ZZ}/\Homology{A}{\ZZ} \to \Torsion{A}{N}(\CC)$.
\end{lemma}

\begin{proof}
Uniformisation (\cref{lem:uniformisation}): the kernel of $\ExpMap{A}$ is
$\PeriodInclusion{A}(\Homology{A}{\ZZ})$, and a point $\ExpMap{A}(x)$ is
torsion if and only if $nx \in \PeriodInclusion{A}(\Homology{A}{\ZZ})$ for
some $n \geq 1$. Then use the torsion of
$\DivisibleGroup{\Homology{A}{\ZZ}}$ (\cref{lem:divisible-torsion}).
Compare \cite[IV.5.2]{DeligneRapoport1973}.
\end{proof}

\begin{corollary}[Tate modules are completed homology]\label{cor:torsion-homology}
\slogan{$A[N] \cong H_1(A,\ZZ)/N$ via division by $N$, and
$\widehat{T}(A) \cong H_1(A,\widehat{\ZZ})$.}
In \cref{setup:abelian}, for $N \geq 1$ the map
$v \mapsto \ExpMap{A}(\PeriodInclusion{A}(v)/N)$ is an isomorphism
$\Homology{A}{\ZZ}/N\Homology{A}{\ZZ} \to \Torsion{A}{N}(\CC)$; it is the
division map $\DivisionMap{N}$ followed by the description of torsion
points (\cref{lem:torsion-points}). Consequently
$\TateModule{A} := \varprojlim_N \Torsion{A}{N}(\CC) \cong \Homology{A}{\ProfiniteIntegers}$,
$\RationalTateModule{A} := \TateModule{A}\otimes\QQ \cong
\Homology{A}{\FiniteAdeles{\QQ}}$ and
$\RationalTateModule{A}/\TateModule{A} \cong A(\CC)_{\mathrm{tors}}$.
\end{corollary}

\begin{proof}
Combine the description of torsion points (\cref{lem:torsion-points}),
which is compatible with multiplication maps, with the Tate modules of
$(\FreeModule\otimes\QQ)/\FreeModule$ (\cref{cor:abstract-tate}).
\end{proof}

\begin{remark}
The intrinsic description of $\Torsion{A}{N}$ is the subgroup
$\tfrac{1}{N}\Homology{A}{\ZZ}/\Homology{A}{\ZZ}$ of
$\Homology{A}{\QQ}/\Homology{A}{\ZZ}$, adapted to the direct system of the
$\Torsion{A}{N}$ under inclusions. The description $\Homology{A}{\ZZ}/N$
uses division by $N$ and is adapted to the inverse system under
multiplication maps, i.e.\ to the Tate module. The factor $1/N$ in the
standard level structure $\StdLevel{\tau}{N}$ is this division by $N$
(\cref{prop:level-from-frames}).
\end{remark}

\begin{lemma}[Weil pairing and Riemann form]\label{lem:weil-riemann}
\slogan{The Weil pairing is the Riemann form modulo $N$.}
In \cref{setup:abelian} with $g = 1$, normalise the Weil pairing as in
\cite{DiamondShurman2005}, so that $\WeilPairing{N}(\tau/N, 1/N) =
\exp(2\pi i/N)$ on $\CC/(\ZZ+\ZZ\tau)$ for $\tau \in \UpperHalfPlane$. Then
under the description of torsion points (\cref{lem:torsion-points}),
$\WeilPairing{N}(x, y) = \exp(2\pi i N\,\RiemannForm{A}(x,y))$ for
$x, y \in \tfrac{1}{N}\Homology{A}{\ZZ}/\Homology{A}{\ZZ}$, where
$\RiemannForm{A}$ is extended $\QQ$-bilinearly. Equivalently, under the
isomorphism $\Homology{A}{\ZZ}/N \cong \Torsion{A}{N}(\CC)$ given by
division by $N$ (\cref{cor:torsion-homology}),
$\WeilPairing{N}(x, y) = \exp(2\pi i\,\RiemannForm{A}(x,y)/N)$ for
$x, y \in \Homology{A}{\ZZ}/N$.
\end{lemma}

\begin{proof}
The right hand sides are well defined since $N\RiemannForm{A}(x,y) \in \ZZ$
for $x \in \Homology{A}{\ZZ}$, $y \in \frac{1}{N}\Homology{A}{\ZZ}$. Both
sides are bilinear and alternating, so it suffices to compare them on one
basis. Choose $z$ with $z(\PeriodInclusion{A}(\Homology{A}{\ZZ})) =
\ZZ+\ZZ\tau$, $\tau \in \UpperHalfPlane$, and the elements of
$\tfrac{1}{N}\Homology{A}{\ZZ}$ mapping to $\tau/N, 1/N$; by the formula
for the Riemann form of an elliptic curve (\cref{lem:elliptic-riemann-form})
$\RiemannForm{A}$ takes the value $\mathrm{Im}(\tau)/(N^{2}\mathrm{Im}(\tau)) = N^{-2}$
on them.
\end{proof}

\section{Frames of elliptic curves}

\begin{definition}[Frames]\label{def:frames}
In \cref{setup:abelian}, a \emph{frame} of $A$ is a group isomorphism
$\Frame \colon \StdModule{\ZZ}{2g} \to \Homology{A}{\ZZ}$, where
$\StdModule{\ZZ}{2g}$ consists of column vectors (\cref{conv:columns}).
Write $\Complexification{\Frame}$ for its $\CC$-linear extension. The group
$\GL{2g}{\ZZ}$ acts on the set of frames on the right by
$\Frame \mapsto \Frame\circ\gamma$. Let
$\StdSymplectic{g} := \smallMat{0}{\IdentityBlock{g}}{-\IdentityBlock{g}}{0}$.
For a Riemann form $\RiemannForm{}$, the frame is
\emph{$\RiemannForm{}$-symplectic} if
$\RiemannForm{}(\Frame v, \Frame w) = \transpose{v}\StdSymplectic{g}w$ for all
$v, w$. If $g = 1$, a frame is \emph{oriented} if it is
$\RiemannForm{A}$-symplectic for the principal Riemann form of
\cref{lem:elliptic-riemann-form}.
\end{definition}

Note that $\StdSymplectic{1} = \Sform$ is the matrix of \cref{lem:S-identity}
and $\transpose{v}\StdSymplectic{1}w = \DetPair{v}{w}$.

\begin{definition}[Period point of a frame]\label{def:period-point}
In \cref{setup:abelian} with $g = 1$, the \emph{period point} of a frame
$\Frame$ is the unique $\PeriodPoint{\Frame} \in \SlitPlane$ with
$\Complexification{\Frame}^{-1}(\HodgeFiltration{A}) =
\CC\,\PeriodVector{\PeriodPoint{\Frame}}$.
\end{definition}

\begin{lemma}[Period points exist]\label{lem:period-point-exists}
\cref{def:period-point} makes sense.
\end{lemma}

\begin{proof}
$\Complexification{\Frame}^{-1}(\HodgeFiltration{A})$ is a line in
$\StdModule{\CC}{2}$ containing no nonzero real vector
(\cref{lem:hodge-kernel}). If it is spanned by $\transpose{(p,q)}$ then
$q \neq 0$ and $p/q \notin \RR$, since otherwise $\transpose{(1,0)}$
respectively $\transpose{(p/q,1)}$ would be a real vector in it. Put
$\PeriodPoint{\Frame} := p/q$.
\end{proof}

\begin{proposition}[Frames and points of the slit plane]\label{prop:frame-tau}
\slogan{An elliptic curve with a frame of its homology is a point of
$\CC \setminus \RR$; the period vector spans the Hodge filtration, and the
lattice map is the period map of a normalised differential.}
\begin{enumerate}
\item Let $\tau \in \SlitPlane$ and identify
  $\Homology{\EllCurve{\tau}}{\ZZ} = \LatticeOf{\tau}$ and
  $\LieAlg{\EllCurve{\tau}} = \CC$ via the uniformisation of
  $\EllCurve{\tau} = \CC/\LatticeOf{\tau}$ (\cref{def:lattice-of-tau}). Then
  $\LatticeBasis{\tau}$ is a frame of $\EllCurve{\tau}$ with period point
  $\tau$, and
  $\RiemannForm{\EllCurve{\tau}}(\LatticeBasis{\tau}v, \LatticeBasis{\tau}w) =
  \mathrm{sign}(\mathrm{Im}\,\tau)\DetPair{v}{w}$. In particular
  $\LatticeBasis{\tau}$ is oriented if and only if $\tau \in \UpperHalfPlane$.
\item The map $(E, \Frame) \mapsto \PeriodPoint{\Frame}$ is a bijection from
  the set of isomorphism classes of pairs of a complex elliptic curve $E$
  and a frame $\Frame$ of $E$ onto $\SlitPlane$, with inverse
  $\tau \mapsto (\EllCurve{\tau}, \LatticeBasis{\tau})$. It restricts to a
  bijection from oriented frames onto $\UpperHalfPlane$.
\item For a frame $\Frame$ and $\gamma \in \GL{2}{\ZZ}$ one has
  $\PeriodPoint{\Frame\circ\gamma} = \MobiusAction{\gamma^{-1}}{\PeriodPoint{\Frame}}$.
\item Let $\Frame$ be an oriented frame of $E$ and $\tau := \PeriodPoint{\Frame}$,
  and let $\FrameDifferential{\Frame} \in \Differentials{E}$ be the
  differential with
  $\PoincareDual{\RiemannForm{E}}{\FrameDifferential{\Frame}} =
  \Complexification{\Frame}(\PeriodVector{\tau})$
  (\cref{prop:poincare-hodge}). Then
  \[
    \Period{\Frame(v)}{\FrameDifferential{\Frame}} = \DetPair{\PeriodVector{\tau}}{v}
    = \LatticeBasis{\tau}(v) \qquad\text{for } v \in \StdModule{\ZZ}{2};
  \]
  in particular $\Period{\Frame(\StdBasis{1})}{\FrameDifferential{\Frame}} = -1$
  and $\Period{\Frame(\StdBasis{2})}{\FrameDifferential{\Frame}} = \tau$. For
  $(E, \Frame) = (\EllCurve{\tau}, \LatticeBasis{\tau})$ the differential
  $\FrameDifferential{\Frame}$ is $dz$.
\end{enumerate}
\end{proposition}

\begin{proof}
(1) By the coordinates of the lattice of $\tau$ (\cref{lem:lattice-coordinates})
$\LatticeBasis{\tau}$ is a frame, and the $\CC$-linear extension of
$\PeriodInclusion{\EllCurve{\tau}}\circ\LatticeBasis{\tau}$ is
$v \mapsto \DetPair{\PeriodVector{\tau}}{v}$, whose kernel is
$\CC\PeriodVector{\tau}$; this kernel is
$\Complexification{\LatticeBasis{\tau}}^{-1}(\HodgeFiltration{\EllCurve{\tau}})$
by the Hodge filtration as kernel of the period map (\cref{lem:hodge-kernel}).
For the orientation, the formula for the Riemann form of an elliptic curve
(\cref{lem:elliptic-riemann-form}) with $z = \mathrm{id}$ gives
$\RiemannForm{\EllCurve{\tau}}(-1, \tau) = \mathrm{Im}(-\overline{\tau})/|\mathrm{Im}\,\tau|
= \mathrm{sign}(\mathrm{Im}\,\tau)$, and an alternating form on
$\StdModule{\ZZ}{2}$ is determined by its value on $(\StdBasis{1},\StdBasis{2})$.

(2) Surjectivity is (1). Let $(E,\Frame)$ and $(E', \Frame')$ have the same
period point. Then $\phi := \Frame'\circ\Frame^{-1}$ is an isomorphism
$\Homology{E}{\ZZ} \to \Homology{E'}{\ZZ}$ whose complexification maps
$\HodgeFiltration{E}$ onto $\HodgeFiltration{E'}$. By \cref{lem:hodge-kernel}
it induces a $\CC$-linear isomorphism
$\LieAlg{E} = \Homology{E}{\CC}/\HodgeFiltration{E} \to
\Homology{E'}{\CC}/\HodgeFiltration{E'} = \LieAlg{E'}$ compatible with the
period inclusions, hence by uniformisation (\cref{lem:uniformisation}) an
isomorphism $E \to E'$ inducing $\phi$ on homology, i.e.\ carrying $\Frame$
to $\Frame'$. The statement about oriented frames follows from (1).

(3) $\Complexification{(\Frame\circ\gamma)}^{-1}(\HodgeFiltration{A}) =
\gamma^{-1}\CC\PeriodVector{\PeriodPoint{\Frame}} =
\CC\PeriodVector{\gamma^{-1}\PeriodPoint{\Frame}}$ by the M\"obius
transformations as left action on lines (\cref{lem:mobius}).

(4) By Poincar\'e duality for the Hodge filtration
(\cref{prop:poincare-hodge}) and since $\Frame$ is oriented,
$\Period{\Frame(v)}{\FrameDifferential{\Frame}} =
\RiemannForm{E}(\Frame\PeriodVector{\tau}, \Frame v) =
\DetPair{\PeriodVector{\tau}}{v}$. For $(\EllCurve{\tau}, \LatticeBasis{\tau})$
the canonical period map (\cref{prop:canonical-period-map}) gives
$\Period{\LatticeBasis{\tau}(v)}{dz} = \LatticeBasis{\tau}(v) =
\DetPair{\PeriodVector{\tau}}{v}$, so $dz$ has the same periods as
$\FrameDifferential{\Frame}$ and the two are equal.
\end{proof}

\begin{remark}[Where the $-1$ comes from]\label{rem:minus-one}
For an oriented frame $\Frame$ with period point $\tau$, the normalised
differential $\FrameDifferential{\Frame}$ has periods $-1$ and $\tau$ along
$\Frame(\StdBasis{1})$ and $\Frame(\StdBasis{2})$, and its row of periods
$(-1, \tau) = \transpose{\PeriodVector{\tau}}\Sform$ is its coordinate vector
in the dual frame of $\Cohomology{E}{\ZZ}$ (see \cref{rem:homology-cohomology}).
The sign cannot be removed while keeping the other conventions: the frame
$\StdBasis{1} \mapsto 1$, $\StdBasis{2} \mapsto \tau$ of
$\Homology{\EllCurve{\tau}}{\ZZ}$ is $\LatticeBasis{\tau}\circ\Diagonal{-1}{1}$,
whose period point is $-\tau$ by (3) of \cref{prop:frame-tau}, which is not
oriented. Classically one normalises a differential by
$\Period{a}{\omega} = 1$ for a basis $(a, b)$ with intersection number
$\IntersectionPairing{a}{b} = 1$ and puts $\tau := \Period{b}{\omega}$;
by the sign convention for Riemann forms (\cref{rem:riemann-sign}) such a
basis is $(-\Frame(\StdBasis{1}), \Frame(\StdBasis{2}))$ for an oriented
frame $\Frame$, and then $\omega = \FrameDifferential{\Frame}$ and the
classical $\tau$ is $\PeriodPoint{\Frame}$. So the $-1$ is the sign
difference between the Riemann form and the intersection pairing.
\end{remark}

\begin{remark}[Homology versus cohomology]\label{rem:homology-cohomology}
The \emph{dual frame} of a frame $\Frame$ is the frame
$\DualFrame{\Frame}$ of $\Cohomology{A}{\ZZ}$ whose coordinates are the
coordinates with respect to the dual basis; changing $\Frame$ to
$\Frame\circ\gamma$ changes $\DualFrame{\Frame}$ to
$\DualFrame{\Frame}\circ\invtranspose{\gamma}$. The map
$\PoincareMap{\RiemannForm{}} \colon \Homology{A}{\ZZ} \to \Cohomology{A}{\ZZ}$,
$v \mapsto \RiemannForm{}(v, \cdot)$, has matrix $\transpose{\StdSymplectic{g}}$
with respect to $\Frame$ and $\DualFrame{\Frame}$ when $\Frame$ is
$\RiemannForm{}$-symplectic. The identity
$\Sform\gamma\Sform^{-1} = \det(\gamma)\invtranspose{\gamma}$ of
\cref{lem:S-identity} is the statement that Poincar\'e duality intertwines
the action on homology frames with the contragredient action on cohomology
frames, up to the similitude factor. Describing $\tau$ through the line
$\HodgeOneZero{E} \subseteq \Cohomology{E}{\CC}$ in the dual frame (the ``dot
product'' convention, spanned by $(-1, \tau)$) instead of through
$\HodgeFiltration{E} \subseteq \Homology{E}{\CC}$ in the frame itself
(spanned by $\transpose{(\tau, 1)}$) is what produces the transposes in the
literature (\cref{rem:literature}).
\end{remark}

\begin{proposition}[Level structures are reductions of frames]
\label{prop:level-from-frames}
\slogan{The standard level structure $\beta_{\tau,N}$ is the frame
$b_{\tau}$ modulo $N$, divided by $N$.}
Let $\tau \in \SlitPlane$ and $N \geq 1$. Then $\StdLevel{\tau}{N}$ is the
reduction of $\LatticeBasis{\tau}$ modulo $N$ followed by division by $N$,
i.e.\ $\StdLevel{\tau}{N}(v) = \LatticeBasis{\tau}(v)/N$ in
$\tfrac{1}{N}\Homology{\EllCurve{\tau}}{\ZZ}/\Homology{\EllCurve{\tau}}{\ZZ} =
\Torsion{\EllCurve{\tau}}{N}$ (\cref{lem:torsion-points,cor:torsion-homology}). Consequently, in the complex
uniformisation (\cref{prop:complex-uniformisation}) the pair $(\tau, g)$
corresponds to $E = \EllCurve{\tau}$ with the level structure
$(\LatticeBasis{\tau} \bmod N)\circ g$, the right action of
$\GL{2}{\ResidueRing{N}}$ is the reduction of the right action of
$\GL{2}{\ZZ}$ on frames, and the left action of $\gamma \in \GL{2}{\ZZ}$ on
$\SlitPlane$ corresponds to $\Frame \mapsto \Frame\circ\gamma^{-1}$ by (3)
of \cref{prop:frame-tau}. By \cref{lem:weil-riemann}, oriented frames
reduce to level structures with Weil determinant $\exp(2\pi i/N)$
(compare \cref{lem:weil-determinant}).
\end{proposition}

\begin{proof}
By definition $\StdLevel{\tau}{N}(v) = \LatticeBasis{\tau}(v)/N \bmod
\LatticeOf{\tau} = \ExpMap{\EllCurve{\tau}}(\PeriodInclusion{\EllCurve{\tau}}
(\LatticeBasis{\tau}v)/N)$.
\end{proof}

\begin{remark}[Rational frames and the big picture]\label{rem:rational-frames}
A point $(\tau, g)$ of the adelic double coset (\cref{def:adelic-double-coset})
gives the curve $\EllCurve{\tau}$, the rational frame
$\LatticeBasis{\tau}\otimes\QQ \colon \StdModule{\QQ}{2} \to
\Homology{\EllCurve{\tau}}{\QQ}$ and the adelic level structure
$\eta := (\LatticeBasis{\tau}\otimes\FiniteAdeles{\QQ})\circ g \colon
\StdModule{\FiniteAdeles{\QQ}}{2} \to \RationalTateModule{\EllCurve{\tau}}$.
For $\gamma \in \GLplus{2}{\QQ}$, multiplication by
$\det(\gamma)(c\tau+d)^{-1}$ maps $\Homology{\EllCurve{\tau}}{\QQ}$ onto
$\Homology{\EllCurve{\gamma\tau}}{\QQ}$ inside $\CC$ (equivariance of the
determinant pairing, \cref{lem:det-pairing-equivariance}), i.e.\ it is a
quasi-isogeny $\EllCurve{\tau} \to \EllCurve{\gamma\tau}$, and it carries
$\eta$ to the level structure attached to $\gamma\cdot(\tau, g)$. So a point
of $\OpenModCurve{\mathcal{K}}(\CC)$ is an elliptic curve up to
quasi-isogeny together with $\eta$ modulo $\mathcal{K}$. The lattices in
$\Homology{E}{\QQ}$ up to $\unit{\QQ}$ correspond to the curves
$\LieAlg{E}/\PeriodInclusion{E}(L)$ isogenous to $E$; via $\eta$ they are the
vertices of the big picture (\cref{lem:restricted-product}). This is the
intrinsic form of the big picture: it is attached to the isogeny class,
and the frame only serves to identify it with the standard one.
\todo[inline]{Compare with the moduli interpretation of Shimura varieties
in \cite{DeligneVarietes1979,Milne2005}, in particular the direction of
$\eta$ and the side on which $\mathcal{K}$ acts.}
\end{remark}

\section{Principally polarised abelian varieties}

\begin{proposition}[Siegel period point of a frame]\label{prop:siegel-period-point}
\slogan{A symplectic frame of a principally polarised abelian variety gives
a point of Siegel space, and the lattice has basis $-1_g$, $Z$.}
In \cref{setup:abelian}, let $\RiemannForm{}$ be a principal Riemann form
and $\Frame$ a $\RiemannForm{}$-symplectic frame. For $Z \in \MatRing{g}{\CC}$
let $\PeriodBlock{Z} := \smallColVec{Z}{\IdentityBlock{g}} \in
\MatRing{2g\times g}{\CC}$.
\begin{enumerate}
\item There is a unique $\SiegelPoint{\Frame} \in \SiegelSpace{g}$, the space
  of complex symmetric $g\times g$ matrices with positive definite
  imaginary part, such that $\Complexification{\Frame}^{-1}(\HodgeFiltration{A})$
  is the column span of $\PeriodBlock{\SiegelPoint{\Frame}}$.
\item For $\gamma \in \Sp{2g}{\ZZ}$ one has
  $\SiegelPoint{\Frame\circ\gamma} = \gamma^{-1}\cdot\SiegelPoint{\Frame}$,
  where $\smallMat{A'}{B'}{C'}{D'}\cdot Z = (A'Z + B')(C'Z+D')^{-1}$.
\item Let $\NormalizedDifferential{\Frame}{j} \in \Differentials{A}$ be the
  differential with $\PoincareDual{\RiemannForm{}}{\NormalizedDifferential{\Frame}{j}}
  = \Complexification{\Frame}(\PeriodBlock{Z}\StdBasis{j})$, where
  $Z = \SiegelPoint{\Frame}$. These form a basis of $\Differentials{A}$, and
  for $x, y \in \StdModule{\ZZ}{g}$
  \[
    \Bigl(\Period{\Frame\smallColVec{x}{y}}{\NormalizedDifferential{\Frame}{j}}\Bigr)_{j}
    = \transpose{\PeriodBlock{Z}}\StdSymplectic{g}\smallColVec{x}{y} = Zy - x .
  \]
  Hence $A(\CC) \cong \CC^{g}/(\ZZ^{g} + Z\ZZ^{g})$.
\end{enumerate}
For $g = 1$ this recovers (1), (3) and (4) of \cref{prop:frame-tau}, with
$\SiegelPoint{\Frame} = \PeriodPoint{\Frame}$ and $\PeriodBlock{\tau} =
\PeriodVector{\tau}$.
\end{proposition}

\begin{proof}
(1) $L := \Complexification{\Frame}^{-1}(\HodgeFiltration{A})$ is a Lagrangian
for $\transpose{v}\StdSymplectic{g}w$ with $-i\,\transpose{u}\StdSymplectic{g}\overline{u} > 0$
for $u \in L\setminus 0$ (Riemann relations, \cref{lem:riemann-relations}).
If $u = \smallColVec{x}{0} \in L$, then $\transpose{u}\StdSymplectic{g}\overline{u} = 0$,
so $u = 0$; hence the projection of $L$ onto the last $g$ coordinates is an
isomorphism and $L$ is the column span of $\PeriodBlock{Z}$ for a unique
$Z$. One computes $\transpose{\PeriodBlock{Z}}\StdSymplectic{g}\PeriodBlock{Z}
= \transpose{Z} - Z$, so isotropy means $Z$ symmetric, and
$\transpose{\PeriodBlock{Z}}\StdSymplectic{g}\overline{\PeriodBlock{Z}} = 2i\,\mathrm{Im}\,Z$,
so positivity means $\mathrm{Im}\,Z$ positive definite. (2) The column span
of $\gamma^{-1}\PeriodBlock{Z}$ is a positive Lagrangian, and writing
$\gamma^{-1} = \smallMat{A'}{B'}{C'}{D'}$ one has
$\gamma^{-1}\PeriodBlock{Z} = \PeriodBlock{\gamma^{-1}\cdot Z}(C'Z+D')$, where
$C'Z+D'$ is invertible by the argument of (1). (3) As in (4) of
\cref{prop:frame-tau}, using Poincar\'e duality for the Hodge filtration
(\cref{prop:poincare-hodge}); the computation is
$\StdSymplectic{g}\smallColVec{x}{y} = \smallColVec{y}{-x}$ and $Z$ symmetric.
The last statement follows from the canonical period map
(\cref{prop:canonical-period-map}) using the coordinates
$\NormalizedDifferential{\Frame}{j}(0)$ on $\LieAlg{A}$; note
$-\ZZ^{g} = \ZZ^{g}$.
\end{proof}

\begin{lemma}[Equivariance of the Siegel lattice map]\label{lem:siegel-equivariance}
\slogan{The higher dimensional version of the equivariance of the
determinant pairing.}
Let $g \geq 1$, $Z \in \SiegelSpace{g}$, and define
$\SiegelLatticeMap{Z} \colon \StdModule{\RR}{2g} \to \CC^{g}$ by
$\SiegelLatticeMap{Z}(v) := \transpose{\PeriodBlock{Z}}\StdSymplectic{g}v$. Let
$\gamma = \smallMat{A'}{B'}{C'}{D'} \in \GSp{2g}{\RR}$ with
$\transpose{\gamma}\StdSymplectic{g}\gamma = \SimilitudeFactor\StdSymplectic{g}$,
$\SimilitudeFactor > 0$. Then $C'Z+D'$ is invertible, $\gamma\cdot Z \in
\SiegelSpace{g}$, and
\[
  \SiegelLatticeMap{\gamma\cdot Z}(\gamma v) =
  \SimilitudeFactor\,\invtranspose{(C'Z+D')}\SiegelLatticeMap{Z}(v)
  \qquad\text{for all } v .
\]
For $g = 1$, $\SimilitudeFactor = \det\gamma$ and this is the equivariance
of the determinant pairing (\cref{lem:det-pairing-equivariance}).
\end{lemma}

\begin{proof}
The first two statements are standard. From
$\gamma\PeriodBlock{Z} = \PeriodBlock{\gamma\cdot Z}(C'Z+D')$ one gets
$\transpose{\PeriodBlock{\gamma\cdot Z}}\StdSymplectic{g}\gamma v =
\invtranspose{(C'Z+D')}\transpose{\PeriodBlock{Z}}\transpose{\gamma}
\StdSymplectic{g}\gamma v$.
\end{proof}

\begin{remark}
For non-principal polarisations of type $D = \mathrm{diag}(d_1, \ldots, d_g)$
one replaces $\StdSymplectic{g}$ by $\smallMat{0}{D}{-D}{0}$; we do not need
this.
\end{remark}

\section{Dictionary}

The following table summarises where the objects live, for an elliptic
curve $E$ with an oriented frame $\Frame$ and $\tau := \PeriodPoint{\Frame}$.

\medskip
\noindent
\begin{tabular}{L{0.29\textwidth}L{0.33\textwidth}L{0.3\textwidth}}
\textbf{Object} & \textbf{Canonical home} & \textbf{In the frame $\Frame$}\\ \hline
cycle & $\Homology{E}{\ZZ}$ & column vector $v \in \StdModule{\ZZ}{2}$\\
Riemann form $\RiemannForm{E}$ & $\textstyle\bigwedge^{2}\Cohomology{E}{\ZZ}$ & $\DetPair{v}{w} = \transpose{v}\Sform w$\\
Hodge filtration & line in $\Homology{E}{\CC}$ & $\CC\PeriodVector{\tau}$, $\PeriodVector{\tau} = \transpose{(\tau,1)}$\\
normalised differential & $\Differentials{E} = \HodgeOneZero{E}$ & row of periods $(-1, \tau)$ in the dual frame\\
its Poincar\'e dual & $\HodgeFiltration{E}$ & $\PeriodVector{\tau}$\\
period of $\FrameDifferential{\Frame}$ along $\Frame(v)$ & $\CC$ & $\DetPair{\PeriodVector{\tau}}{v} = \LatticeBasis{\tau}(v)$\\
Lie algebra & $\Homology{E}{\CC}/\HodgeFiltration{E}$ & $\CC$ with coordinate $z$, $dz = \FrameDifferential{\Frame}$\\
lattice & $\PeriodInclusion{E}(\Homology{E}{\ZZ}) \subseteq \LieAlg{E}$ & $\LatticeOf{\tau} = \ZZ + \ZZ\tau$\\
torsion points & $\Homology{E}{\QQ}/\Homology{E}{\ZZ}$ & $\StdModule{\QQ}{2}/\StdModule{\ZZ}{2}$\\
$N$-torsion & $\tfrac{1}{N}\Homology{E}{\ZZ}/\Homology{E}{\ZZ}$ & $\tfrac{1}{N}\StdModule{\ZZ}{2}/\StdModule{\ZZ}{2}$; $\StdLevel{\tau}{N} = (\Frame \bmod N)/N$\\
Weil pairing & on $\tfrac{1}{N}\Homology{E}{\ZZ}/\Homology{E}{\ZZ}$ & $\exp(2\pi i N\DetPair{v}{w})$\\
rational Tate module & $\Homology{E}{\FiniteAdeles{\QQ}}$ & $\StdModule{\FiniteAdeles{\QQ}}{2}$\\
big picture & lattices in $\Homology{E}{\QQ}$ mod $\unit{\QQ}$ & lattices in $\StdModule{\QQ}{2}$ mod $\unit{\QQ}$\\
change of frame & $\Frame \mapsto \Frame\circ\gamma$ & $\tau \mapsto \gamma^{-1}\tau$
\end{tabular}
\medskip
